% Source PDF pages 81--84; manual pages 2-52--2-55.
\subsection{Atmosphere Model}\label{sec:atmosphere-model}

The atmosphere models available in TAOS were originally developed by Acurex
Corporation under contract to Sandia and subsequently corrected by Landrum in
1987.\manualref{12} The models include the 1976 U.S. standard
atmosphere,\manualref{13} the 1966 supplemental atmospheres,\manualref{14}
Tonopah range data, and a mean annual Kwajalein atmosphere.\manualref{15}

The 1976 U.S. standard atmosphere is an idealized mid-latitude representation of
the atmosphere under year-round moderate solar activity. The other atmosphere
models likewise represent average conditions for a location and season. Actual
atmospheric properties on a particular day can differ substantially from the model;
density differences as large as 50 percent are possible. Because future atmospheric
conditions at a specified location and date cannot be predicted exactly, model
atmospheres are used as a best approximation. The 1976 U.S. standard atmosphere
is widely used to provide consistency when trajectories are compared.

All available models use the same simplifying assumptions. The resulting equations
relate altitude, temperature, pressure, density, speed of sound, and viscosity and
produce the following TAOS output variables:
\begin{center}
\begin{tabularx}{0.78\linewidth}{@{}>{\ttfamily}l>{$}l<{$}X@{}}
  temp   & T    & temperature.\\
  pres   & p    & pressure.\\
  rho    & \rho & density.\\
  sndspd & c    & speed of sound.\\
  nu     & \nu  & kinematic viscosity.\\
\end{tabularx}
\end{center}

\taossubhead{Assumptions}

The atmospheric gas composition is assumed to satisfy the aerostatic equation
\begin{equation}
  dp=-\rho g\,dh
  \label{eq:atmos-aerostatic}
\end{equation}
and the perfect-gas law
\begin{equation}
  p=\rho\frac{R^*}{M_w}T,
  \label{eq:atmos-perfect-gas}
\end{equation}
where $g$ is acceleration due to gravity, $h$ is altitude above sea level, $R^*$ is
the universal gas constant, $8314.32\ \mathrm{N\,m/(kmol\,K)}$, and $M_w$ is
molecular weight.

The atmosphere is divided into concentric layers surrounding the earth. Within each
layer the gradient of molecular temperature $\tau$ with respect to geopotential
altitude $H$ is assumed constant:
\begin{equation}
  \frac{d\tau}{dH}=\text{constant}.
  \label{eq:molecular-temperature-gradient}
\end{equation}
Molecular temperature and geopotential altitude are defined by
\begin{equation}
  \tau=\frac{M_{w0}}{M_w}T
  \label{eq:molecular-temperature-definition}
\end{equation}
and
\begin{equation}
  H=\frac{1}{G}\int_0^h g\,dh.
  \label{eq:geopotential-altitude-integral}
\end{equation}
The zero subscript denotes sea-level conditions. For example, $M_{w0}$ is the
sea-level molecular weight: $28.9644\ \mathrm{kg/kmol}$ for the 1976 U.S.
standard atmosphere and $28.96\ \mathrm{kg/kmol}$ for the other TAOS atmosphere
models. The reference geopotential is
$G=9.80665\ \mathrm{m/s^2}$, corresponding to the sea-level gravitational
acceleration at latitude $45.5425^\circ$ from
\cref{eq:lambert-sea-level-gravity}.

For atmospheric calculations, gravity is assumed to vary with altitude according to
the inverse-square law
\begin{equation}
  g=g_0\left(\frac{R'_{\oplus}}{R'_{\oplus}+h}\right)^2,
  \label{eq:atmos-gravity-inverse-square}
\end{equation}
where $R'_{\oplus}$ is an effective earth radius used only for geopotential
calculations. Substitution into \cref{eq:geopotential-altitude-integral} and
integration gives
\begin{equation}
  H=\frac{g_0}{G}\left(\frac{R'_{\oplus}h}{R'_{\oplus}+h}\right).
  \label{eq:geopotential-altitude-closed-form}
\end{equation}
The sea-level gravitational acceleration, in $\mathrm{m/s^2}$, is assumed to vary
with latitude according to Lambert's equation:
\begin{equation}
  g_0=9.806160\left(1-0.0026373\cos 2\delta
                      +0.0000059\cos^2 2\delta\right).
  \label{eq:lambert-sea-level-gravity}
\end{equation}
Here $\delta$ is the latitude associated with the atmosphere model, not necessarily
the current trajectory latitude. For example, a $60^\circ$N atmosphere uses
$\delta=60^\circ$. This convention is inconsistent with the gravity model used for
trajectory propagation, but it is consistent with the 1976 U.S. standard and 1966
supplemental atmosphere definitions.

The effective earth radius in \cref{eq:geopotential-altitude-closed-form} is
\begin{equation}
  R'_{\oplus}
  =\frac{2g_0}
  {3.085462\times10^{-6}
   +2.27\times10^{-9}\cos 2\delta
   -2.0\times10^{-12}\cos 4\delta}.
  \label{eq:effective-geopotential-earth-radius}
\end{equation}
At latitude $45.5425^\circ$, this expression gives
$R'_{\oplus}=6356.766\ \mathrm{km}$.

The properties at the base of each atmosphere layer are supplied in tables as
functions of geometric altitude $h$ or geopotential altitude $H$. Molecular
temperature $\tau$ and its gradient $d\tau/dH$ are tabulated against $H$, while
molecular weight $M_w$ is tabulated against $h$.

\taossubhead{Temperature}

For a point within the $i$th layer, molecular temperature and molecular weight are
computed from the base values as
\begin{equation}
  \tau=\tau_i+\frac{d\tau}{dH}\left(H-H_i\right)
  \label{eq:layer-molecular-temperature}
\end{equation}
and
\begin{equation}
  M_w=M_{wi}+\frac{dM_w}{dH}\left(H-H_i\right).
  \label{eq:layer-molecular-weight}
\end{equation}
The atmospheric temperature is then
\begin{equation}
  T=\frac{M_w}{M_{w0}}\tau.
  \label{eq:atmos-temperature}
\end{equation}

\taossubhead{Pressure}

Writing \cref{eq:atmos-aerostatic,eq:atmos-perfect-gas} in terms of geopotential
altitude using \cref{eq:molecular-temperature-definition,eq:geopotential-altitude-integral}
gives
\begin{equation}
  dp=-\rho g_0\,dH
  \label{eq:pressure-geopotential-differential}
\end{equation}
and
\begin{equation}
  p=\frac{\rho R^*\tau}{M_{w0}}.
  \label{eq:pressure-molecular-temperature}
\end{equation}
Dividing the first expression by the second gives
\begin{equation}
  \frac{dp}{p}
  =-\frac{g_0M_{w0}}{R^*\tau}\,dH.
  \label{eq:pressure-separable-differential}
\end{equation}
Because $\tau$ is linear in $H$ within a layer, integration gives one of two forms.
For an isothermal layer,
\begin{equation}
  p=p_i\exp\left[-\frac{g_0M_{w0}(H-H_i)}{R^*\tau_i}\right],
  \qquad \frac{d\tau}{dH}=0,
  \label{eq:pressure-isothermal-layer}
\end{equation}
while for a nonzero temperature gradient,
\begin{equation}
  p=p_i
  \left(\frac{\tau_i}{\tau}\right)^{
    \frac{g_0M_{w0}}{R^*\left(d\tau/dH\right)}},
  \qquad \frac{d\tau}{dH}\ne0.
  \label{eq:pressure-gradient-layer}
\end{equation}
These formulas are first used to determine pressure at atmosphere-layer base
altitudes and then to compute pressure at an arbitrary geopotential altitude.

\taossubhead{Density, Speed of Sound, and Viscosity}

Combining \cref{eq:atmos-perfect-gas,eq:molecular-temperature-definition} gives
atmospheric density:
\begin{equation}
  \rho=\frac{pM_{w0}}{R^*\tau}.
  \label{eq:atmos-density}
\end{equation}
The speed of sound is
\begin{equation}
  c=\sqrt{\frac{kR^*T}{M_w}}
   =\sqrt{\frac{kR^*\tau}{M_{w0}}},
  \label{eq:atmos-speed-of-sound}
\end{equation}
where $k$ is the ratio of specific heats, taken as $1.4$ for air. Kinematic
viscosity is computed from Sutherland's formula:
\begin{equation}
  \nu
  =\frac{1}{\rho}
   \frac{1.458\times10^{-6}\,\tau^{1.5}}{\tau+110.4}.
  \label{eq:atmos-kinematic-viscosity}
\end{equation}

\taossubhead{Atmosphere Calculations}

The 1976 U.S. standard-atmosphere table is defined through an altitude of
$146\ \mathrm{km}$. The other models are defined to altitudes between
$30\ \mathrm{km}$ and $117\ \mathrm{km}$ and are extrapolated so that they meet
the 1976 standard atmosphere at $146\ \mathrm{km}$. Before trajectory
calculations, the function \texttt{atmos\_setup} selects the requested atmosphere
model and loads the applicable tables.

Below $146\ \mathrm{km}$, the function \texttt{atmos\_properties} evaluates the
model using the preceding equations. Above $146\ \mathrm{km}$, all models use a
separate 1976-standard-atmosphere table containing temperature, pressure, and
density through $1000\ \mathrm{km}$. The function \texttt{atmos\_hi\_alt}
interpolates temperature linearly and pressure and density exponentially. All model
atmosphere functions reside in the \texttt{atmos} module so they can be reused by
other software.
